% !TEX root = ../main.tex
\begin{ouverture}[interferences]{22}{Diffraction et interférences}{Partie 4 · Ondes et signaux}
{Diffraction\enspace·\enspace Angle caractéristique\enspace·\enspace Interférences constructives et destructives\enspace·\enspace Fentes d'Young}
\objectif{Caractériser le phénomène de diffraction dans des situations variées et en citer des conséquences concrètes.}
\objectif{Exploiter la relation entre l'angle caractéristique de diffraction, la longueur d'onde et la taille de
l'ouverture.}
\objectif{Établir les conditions d'interférences constructives et destructives de deux ondes issues de sources en phase.}
\objectif{Prévoir les franges d'interférences des fentes d'Young et établir l'expression de l'interfrange.}
\objectif{Représenter en Python la somme de deux signaux sinusoïdaux synchrones de déphasage variable.}
\prerequis{Onde progressive périodique, période, longueur d'onde, célérité : $\lambda=v\,T$ (classe de première).}
\prerequis{Lumière, spectre, ondes électromagnétiques (classes de seconde et de première).}
\prerequis{Fonctions trigonométriques (mathématiques).}
\end{ouverture}
\essentiel
\essSection{1}{La diffraction}
\begin{definition}[Diffraction]
Lorsqu'une onde rencontre une ouverture ou un obstacle, elle se propage au-delà dans des directions qu'elle n'aurait pas
prises en ligne droite : c'est la \textbf{diffraction}. Elle est d'autant plus marquée que la taille $a$ de l'ouverture
ou de l'obstacle est petite devant la longueur d'onde $\lambda$ ; pour $a\gg\lambda$, elle devient négligeable. La
diffraction ne modifie ni la fréquence ni la longueur d'onde.
\end{definition}
\begin{loi}[Angle caractéristique de diffraction]
Pour une fente de largeur $a$, la tache centrale de diffraction est vue sous le demi-angle
\[\theta=\frac{\lambda}{a},\]
avec $\theta$ en radians ; cette relation est valable pour les petits angles. Sur un écran à la distance $D$, la tache
centrale a pour largeur $L=2D\tan\theta\approx\dfrac{2\lambda D}{a}$.
\end{loi}
\schema{22-diffraction}
\begin{exemple}
Un laser rouge ($\lambda=\qty{633}{nm}$) éclaire une fente de \qty{0,10}{mm} : $\theta=\qty{6,3e-3}{rad}$. Sur un écran à
\qty{2,0}{m}, la tache centrale mesure \qty{2,5}{cm}, alors que le faisceau n'avait que quelques millimètres. La
diffraction explique qu'on entende un son derrière une porte ouverte ($\lambda\approx\qty{1}{m}$), que la houle contourne
les jetées, et qu'elle limite la finesse des images de tout instrument d'optique.
\end{exemple}
\essSection{2}{Interférences de deux ondes}
\begin{definition}[Interférences]
Deux ondes de même fréquence, dont le déphasage est constant (on dit qu'elles sont \textbf{cohérentes}), se superposent :
en certains points, leurs effets s'ajoutent, l'amplitude est maximale (interférences \textbf{constructives}) ; en
d'autres, ils se compensent, l'amplitude est minimale, voire nulle (interférences \textbf{destructives}).
\end{definition}
\schema{22-sommes}
\begin{propriete}[Conditions d'interférences]
Deux sources ponctuelles $S_1$ et $S_2$ émettent en phase des ondes de longueur d'onde $\lambda$ dans un milieu
homogène. En un point $M$, on note $\delta=S_2M-S_1M$ la \textbf{différence de marche}. Les interférences sont :
\begin{itemize}
\item constructives si $\delta=k\,\lambda$, avec $k$ entier relatif ;
\item destructives si $\delta=\left(k+\dfrac12\right)\lambda$.
\end{itemize}
\end{propriete}
\begin{savoirfaire}[Établir les conditions d'interférences]
L'onde issue de $S_2$ parcourt $\delta$ de plus que celle issue de $S_1$ : elle arrive en $M$ avec le retard
$\tau=\dfrac{\delta}{v}$. Les deux ondes sont en phase si ce retard est un nombre entier de périodes, $\tau=k\,T$, soit
$\delta=k\,v\,T=k\,\lambda$ ; elles sont en opposition de phase si $\tau=\left(k+\frac12\right)T$, soit
$\delta=\left(k+\frac12\right)\lambda$.
\end{savoirfaire}
\essSection{3}{Interférences lumineuses : les fentes d'Young}
\begin{propriete}[Franges d'interférences]
Deux fentes fines $S_1$ et $S_2$, distantes de $b$ et éclairées par le même laser, se comportent comme deux sources
cohérentes en phase. Sur un écran à la distance $D\gg b$, au point d'abscisse $x$ mesurée depuis le centre, la
différence de marche vaut $\delta\approx\dfrac{b\,x}{D}$ (expression fournie). Les franges brillantes ($\delta=k\lambda$)
sont situées en $x_k=k\,\dfrac{\lambda D}{b}$ : elles sont équidistantes, séparées par l'\textbf{interfrange}
\[i=\frac{\lambda\,D}{b}.\]
\end{propriete}
\schema{22-young}
\begin{exemple}
Avec $\lambda=\qty{633}{nm}$, $b=\qty{0,20}{mm}$ et $D=\qty{1,50}{m}$ : $i=\qty{4,75}{mm}$. Mesurer l'interfrange permet donc de
déterminer une longueur d'onde. Les interférences expliquent aussi les irisations des bulles de savon, les traitements
antireflets et le fonctionnement des détecteurs d'ondes gravitationnelles.
\end{exemple}
\begin{codecote}[0.64]
\begin{pythoncol}
from math import sin, pi
A, T = 1.0, 1.0
instants = [k / 1000 for k in range(1000)]
for phi in [0, pi/2, 2*pi/3, pi]:
    s = [A*sin(2*pi*t/T) + A*sin(2*pi*t/T + phi)
         for t in instants]
    print(round(phi, 2), round(max(s), 2))
\end{pythoncol}
\begin{consolecol}
0 2.0
1.57 1.41
2.09 1.0
3.14 0.0
\end{consolecol}
\cote
\fontsize{9.5}{12.5}\selectfont
\capacite{Capacité numérique}\par\smallskip
Le programme additionne, sur une période, deux signaux sinusoïdaux de même amplitude $A$ et de même période, le
second étant déphasé de $\varphi$, puis affiche l'amplitude de la somme : $2A$ en phase, $\sqrt2\,A$ pour
$\varphi=\pi/2$, $A$ pour $\varphi=2\pi/3$, zéro en opposition de phase.
\end{codecote}
\begin{attention}
Dans $\theta=\lambda/a$, l'angle s'exprime en radians et $a$ désigne la largeur de la fente ; dans $i=\lambda D/b$, $b$ est la
distance entre les deux fentes. Deux lampes distinctes n'interfèrent pas : leurs ondes ne sont pas cohérentes.
\end{attention}
\begin{experience}[Interférences à la surface de l'eau]
Dans une cuve à ondes, deux pointes reliées au même vibreur frappent la surface en phase. On observe des lignes où
l'eau reste immobile, de forme hyperbolique, séparées par des zones de forte agitation. En relevant les distances
$S_1M$ et $S_2M$ en des points immobiles, on vérifie que $\delta$ vaut $\lambda/2$, $3\lambda/2$\,\dots Avec un laser et des
fentes d'Young, on mesure de même l'interfrange, éventuellement sur une photographie traitée par un logiciel, pour en
déduire la longueur d'onde.
\end{experience}
\begin{rediger}
Pour conclure sur la nature des interférences en un point, calculer la différence de marche, puis le rapport
$\delta/\lambda$ : « $\delta=2{,}5\,\lambda$, demi-entier : interférences destructives ». Pour mesurer un interfrange, mesurer
plusieurs interfranges à la fois afin de réduire l'incertitude.
\end{rediger}
\exercices

\rubrique{Découverte}
\begin{blocLivre}{ex-22-1}
\begin{exercice}[1]{Diffraction ou pas ?}
\begin{questions}
\item On entend une conversation à travers une porte ouverte sans voir les personnes qui parlent. Les longueurs d'onde
des sons de la voix sont comprises entre \qty{0,3}{m} et \qty{3}{m}, celles de la lumière visible entre 400 et \qty{800}{nm}.
Expliquer.
\item Dans quel cas la diffraction d'une onde par une ouverture est-elle marquée ?
\item La diffraction modifie-t-elle la fréquence de l'onde ? sa longueur d'onde ?
\item Citer deux situations de la vie courante où la diffraction se manifeste.
\end{questions}
\end{exercice}
\end{blocLivre}
\begin{blocLivre}{ex-22-2}
\begin{exercice}[1]{Mesurer la largeur d'une fente}
Un laser de longueur d'onde \qty{650}{nm} éclaire une fente. Sur un écran placé à \qty{1,80}{m}, la tache centrale de
diffraction mesure \qty{2,3}{cm}.
\begin{questions}
\item Représenter la situation par un schéma où figurent l'angle $\theta$, la distance $D$ et la largeur $L$.
\item Exprimer $\theta$ en fonction de $L$ et de $D$, en utilisant l'approximation des petits angles, puis le calculer.
\item En déduire la largeur de la fente.
\item Comment évolue la tache centrale si l'on utilise une fente deux fois plus fine ?
\end{questions}
\end{exercice}
\end{blocLivre}
\begin{blocLivre}{ex-22-3}
\begin{exercice}[1]{La houle entre dans le port}
Une houle de longueur d'onde \qty{40}{m} arrive perpendiculairement à l'entrée d'un port, large de \qty{60}{m}.
\begin{questions}
\item La houle est-elle diffractée par l'entrée du port ? Justifier.
\item Calculer l'angle caractéristique de diffraction, en radians puis en degrés. On admettra que la relation
$\theta=\lambda/a$ donne encore un ordre de grandeur.
\item Pourquoi des bateaux amarrés loin de l'axe de l'entrée sont-ils malgré tout agités ?
\item Que deviendrait l'angle pour une passe de \qty{400}{m} ? Commenter.
\end{questions}
\end{exercice}
\end{blocLivre}
\begin{blocLivre}{ex-22-4}
\begin{exercice}[1]{Constructives ou destructives ?}
Deux pointes $S_1$ et $S_2$ frappent en phase la surface de l'eau ; les ondes ont une longueur d'onde de \qty{2,0}{cm}.
\begin{center}
\begin{tabular}{lcccc}
\toprule
Point & $M$ & $N$ & $P$ & $Q$ \\
$S_1M$ (cm) & 8,0 & 6,0 & 9,0 & 5,0 \\
$S_2M$ (cm) & 11,0 & 10,0 & 9,0 & 10,0 \\
\bottomrule
\end{tabular}
\end{center}
\begin{questions}
\item Rappeler les conditions d'interférences constructives et destructives.
\item Calculer, pour chaque point, la différence de marche et le rapport $\delta/\lambda$.
\item En quels points l'eau est-elle immobile ? En quels points l'agitation est-elle maximale ?
\item Où se trouvent tous les points tels que $S_1M=S_2M$ ? Que vaut l'amplitude en ces points ?
\end{questions}
\end{exercice}
\end{blocLivre}
\begin{blocLivre}{ex-22-5}
\begin{exercice}[1]{Les fentes d'Young}
Deux fentes distantes de $b=\qty{0,25}{mm}$ sont éclairées par un laser vert de longueur d'onde \qty{532}{nm}. L'écran est à
$D=\qty{2,00}{m}$.
\begin{questions}
\item Calculer l'interfrange.
\item Que devient l'interfrange si l'on recule l'écran à \qty{4,00}{m} ?
\item Même question avec un laser rouge de \qty{650}{nm}, l'écran étant à \qty{2,00}{m}.
\item Comment varie l'interfrange si l'on rapproche les deux fentes ?
\end{questions}
\end{exercice}
\end{blocLivre}
\rubrique{Entraînement}
\begin{blocLivre}{ex-22-6}
\begin{exercice}[2]{Mesurer un cheveu}
Un cheveu tendu devant un laser de longueur d'onde \qty{633}{nm} produit la même figure de diffraction qu'une fente de même
largeur. Sur un écran à \qty{2,00}{m}, la tache centrale mesure \qty{4,2}{cm}.
\begin{questions}
\item Calculer l'angle caractéristique de diffraction.
\item En déduire le diamètre du cheveu, en micromètres.
\item Le diamètre d'un cheveu humain est compris entre 50 et \qty{100}{\micro m}. Le résultat est-il plausible ?
\item Pourquoi la mesure directe avec une règle serait-elle impossible ? Quel est l'intérêt de la méthode ?
\end{questions}
\end{exercice}
\end{blocLivre}
\begin{blocLivre}{ex-22-7}
\begin{exercice}[2]{Diffraction de la lumière blanche}
Une fente de \qty{50}{\micro m} est éclairée en lumière blanche ; l'écran est à \qty{1,5}{m}.
\begin{questions}
\item Calculer la largeur de la tache centrale pour une lumière bleue de \qty{400}{nm}, puis pour une lumière rouge de
\qty{750}{nm}.
\item Quelle couleur est la plus diffractée ?
\item Décrire l'aspect du centre et des bords de la tache centrale en lumière blanche.
\item Pourquoi utilise-t-on un laser, plutôt qu'une lampe, pour mesurer une longueur d'onde par diffraction ?
\end{questions}
\end{exercice}
\end{blocLivre}
\begin{blocLivre}{ex-22-8}
\begin{exercice}[2]{La somme de deux signaux en Python}
On reprend le programme de l'essentiel.
\begin{questions}
\item \capacite{Capacité numérique}Que contient la liste \code{instants} ? Que représente la liste \code{s} ?
\item Comment le programme détermine-t-il l'amplitude de la somme ? Pourquoi suffit-il d'étudier une seule période ?
\item On admet que l'amplitude de la somme vaut $2A\left|\cos\dfrac{\varphi}{2}\right|$. Vérifier les quatre valeurs affichées.
\item Écrire les lignes qui tracent, sur un même graphique, les deux signaux et leur somme pour $\varphi=\pi/2$.
\end{questions}
\end{exercice}
\end{blocLivre}
\begin{blocLivre}{ex-22-9}
\begin{exercice}[2]{Mesurer une longueur d'onde}
Des fentes d'Young distantes de $b=\qty{0,40}{mm}$ sont éclairées par une diode laser. Sur un écran à $D=\qty{2,50}{m}$, on
mesure \qty{3,8}{cm} entre les centres de deux franges brillantes séparées par dix interfranges.
\begin{questions}
\item Pourquoi mesure-t-on dix interfranges plutôt qu'un seul ?
\item Calculer l'interfrange, puis la longueur d'onde de la diode laser. Quelle est sa couleur ?
\item La mesure sur l'écran est faite à \qty{1}{mm} près. Estimer l'incertitude relative sur l'interfrange.
\item Établir l'expression de l'interfrange à partir de $\delta=\dfrac{b\,x}{D}$.
\end{questions}
\end{exercice}
\end{blocLivre}
\begin{blocLivre}{ex-22-10}
\begin{exercice}[2]{La cuve à ondes}
Dans une cuve à ondes, deux pointes distantes de $S_1S_2=\qty{6,0}{cm}$ vibrent en phase à \qty{20}{Hz}. Les ondes se propagent
à \qty{0,30}{m.s^{-1}}.
\begin{questions}
\item Calculer la longueur d'onde.
\item Sur le segment $[S_1S_2]$, la différence de marche varie entre \qty{-6,0}{cm} et \qty{6,0}{cm}. Combien de lignes de
repos coupent ce segment ? Combien de lignes d'amplitude maximale, en dehors des sources ?
\item Montrer que deux points de repos consécutifs du segment sont séparés de $\lambda/2$.
\item Comment évoluerait la figure si l'on doublait la fréquence ?
\end{questions}
\end{exercice}
\end{blocLivre}
\begin{blocLivre}{ex-22-11}
\begin{exercice}[2]{Deux sources incohérentes}
On éclaire deux fentes d'Young, non plus avec un laser, mais avec deux lampes identiques placées chacune derrière une
fente.
\begin{questions}
\item Observe-t-on des franges ? Justifier.
\item Pourquoi un laser, qui éclaire les deux fentes à la fois, produit-il des franges stables ?
\item Deux haut-parleurs alimentés par le même générateur basse fréquence produisent-ils des interférences ? Et deux
musiciens qui jouent la même note ?
\item En lumière blanche, avec une seule source éclairant les deux fentes, la frange centrale est blanche et les
suivantes sont irisées. Expliquer à l'aide de l'expression de l'interfrange.
\end{questions}
\end{exercice}
\end{blocLivre}
\rubrique{Objectif bac}
\begin{blocLivre}{ex-22-12}
\begin{exercice}[3]{La finesse des images d'un télescope}
\begin{docu}{1}{La limite de diffraction}
L'ouverture circulaire d'un instrument d'optique diffracte la lumière : l'image d'une étoile n'est pas un point mais une
petite tache. Deux points ne sont distingués que s'ils sont vus sous un angle supérieur à $\theta=1{,}22\,\dfrac{\lambda}{D}$,
où $D$ est le diamètre de l'ouverture.
\end{docu}
\begin{docu}{2}{Trois instruments}
Le télescope spatial Hubble a un miroir de \qty{2,4}{m} de diamètre et observe en lumière visible
($\lambda=\qty{550}{nm}$) ; le télescope James-Webb, \qty{6,5}{m}, observe en infrarouge ($\lambda\approx\qty{2}{\micro m}$) ; la pupille
de l'œil mesure environ \qty{3}{mm} en plein jour.
\end{docu}
\begin{donnees}
Distance Terre-Lune : \qty{3,84e8}{m}.
\end{donnees}
\begin{questions}
\item Calculer l'angle limite de résolution de Hubble.
\item Quelle est la taille du plus petit détail de la surface lunaire que Hubble peut distinguer ?
\item \initiative Comparer la résolution de l'œil et celle de Hubble pour l'observation de la Lune.
\item Le James-Webb, plus grand, a-t-il une meilleure résolution que Hubble ? Pourquoi construit-on des télescopes de
plus en plus grands ?
\end{questions}
\end{exercice}
\end{blocLivre}
\begin{blocLivre}{ex-22-13}
\begin{exercice}[3]{Le casque à réduction de bruit active}
\begin{docu}{1}{Combattre le bruit par le bruit}
Un micro capte le bruit extérieur ; l'électronique du casque fait émettre par l'écouteur une onde de même amplitude, en
opposition de phase. Les deux ondes interfèrent de façon destructive dans le conduit auditif.
\end{docu}
\begin{docu}{2}{L'importance du retard}
Si l'onde émise par l'écouteur a un retard $\Delta t$ par rapport à l'opposition de phase parfaite, l'amplitude résiduelle
vaut $2A\,\sin(\pi f\,\Delta t)$, $A$ étant l'amplitude du bruit. Une amplitude divisée par 10 correspond à une atténuation
de \qty{20}{dB}.
\end{docu}
\begin{questions}
\item Pour un bruit de \qty{200}{Hz}, de quelle durée faut-il décaler un signal pour le mettre en opposition de phase avec
lui-même ?
\item Montrer que l'amplitude résiduelle est nulle pour $\Delta t=0$.
\item \initiative Quel retard maximal l'électronique peut-elle tolérer pour atténuer de \qty{20}{dB} au moins un bruit de
\qty{200}{Hz} ? un bruit de \qty{2000}{Hz} ? On admettra que $\sin u\approx u$ pour $u$ petit.
\item Expliquer pourquoi ces casques sont efficaces contre le grondement d'un moteur d'avion, mais beaucoup moins contre
les voix.
\end{questions}
\end{exercice}
\end{blocLivre}
\begin{blocLivre}{ex-22-14}
\begin{exercice}[3]{Le traitement antireflet}
\begin{docu}{1}{Une couche mince}
Les verres de lunettes sont recouverts d'une couche transparente de fluorure de magnésium, d'indice $n=1{,}38$ et
d'épaisseur $e$. La lumière se réfléchit en partie sur la face supérieure de la couche et en partie sur le verre.
\end{docu}
\begin{docu}{2}{La différence de marche}
Les deux ondes réfléchies présentent une différence de marche $\delta=2\,n\,e$ (les deux réflexions introduisent le même
déphasage, qui ne joue donc aucun rôle).
\end{docu}
\begin{questions}
\item Quelle condition sur $\delta$ supprime la lumière réfléchie de longueur d'onde $\lambda$ ?
\item \initiative Calculer la plus petite épaisseur de couche qui supprime la réflexion à \qty{550}{nm}, au centre du spectre
visible.
\item Calculer le rapport $\delta/\lambda$ pour une lumière bleue de \qty{450}{nm} et pour une lumière rouge de \qty{700}{nm}.
Ces couleurs sont-elles éteintes ?
\item Expliquer la teinte violacée des reflets sur les verres traités.
\end{questions}
\end{exercice}
\end{blocLivre}
\begin{blocLivre}{ex-22-15}
\begin{exercice}[3]{Mesurer la vitesse du son par interférences}
\begin{docu}{1}{Le dispositif}
Deux haut-parleurs, face à face et alimentés par le même générateur à \qty{2000}{Hz}, émettent en phase. Un micro se déplace
sur le segment qui les joint ; un oscilloscope affiche l'amplitude du signal reçu.
\end{docu}
\begin{docu}{2}{Les mesures}
On repère dix positions successives où l'amplitude est minimale. Les positions extrêmes sont distantes de
\qty{77,4}{cm}.
\end{docu}
\begin{questions}
\item Quand le micro se déplace de $x$ vers l'un des haut-parleurs, de combien varie la différence de marche ?
\item En déduire que deux minimums successifs sont séparés de $\lambda/2$.
\item \initiative Calculer la longueur d'onde, puis la célérité du son.
\item Pourquoi est-il préférable de repérer dix minimums plutôt que deux ?
\end{questions}
\end{exercice}
\end{blocLivre}
